\section{Choosing both the operations and the weights}\label{app:synthesis}
We asked whether Astra could implement square, sine or positive self-power when it could choose the circuit topology, meaning the operations and their connections. A domain-specific language (DSL) restricted the allowed operations. Each reply was an instruction list written in JavaScript Object Notation (JSON), a text format for structured data. A wire carries the value computed at one graph node. An input-dependent wire can vary with the input, while an input-independent wire carries a fixed constant. The affine/rectifier operator set permits addition, subtraction, multiplication by an input-independent constant, division by a constant and rectification. A second operator set also permits products of two input-dependent wires. In both sets, constants can use sine, logarithms and other elementary functions when the graph is built. For each unseen input, inference executes only the resulting fixed circuit. Every operation that depends on the input counts once toward the 4,096-operation cap after expansion of loops with fixed bounds. Precomputed constants, the size of the instruction text and the number of wires are separate costs.

Each prompt supplied the target formula, input interval, arithmetic rules and operation cap. Eighteen independent requests covered three targets, two operator sets and three replicate labels (9213, 2451 and 8282). The labels distinguish calls with the same instruction, without changing their inputs. The model returned weights and topology without a given Xavier array, tools, examples or evaluation feedback. In the affine/rectifier self-power condition two requests timed out after 300 seconds. We repeated those two before scoring, with a 900-second limit. One reply had invalid JSON syntax and the other exceeded the operation cap. Neither produced an evaluable circuit. All 16 responses returned within the original budget compiled. Table~\ref{tab:circuit-synthesis} keeps all 18 conditions.

\begin{table}[htbp]
\centering\small\setlength{\tabcolsep}{4pt}
\caption{WSP with free topology. For each trial the accuracy columns count exact outputs out of 16,384 outer inputs against correctly rounded MPFR values. Sine and square use both signs with magnitudes from 3 to 10, and self-power uses positive inputs. The last column gives the number of input-dependent operations in trial order. An invalid reply has no numerical prediction or RMSE.}
\label{tab:circuit-synthesis}
\begin{tabular}{@{}llrrrl@{}}\toprule
& & \multicolumn{3}{c}{Exact outer predictions} & Operation counts\\
Target & Operators & Trial 1 & Trial 2 & Trial 3 & Trials 1, 2, 3\\\midrule
Square & Affine/ReLU & 16,384 & 16,384 & 16,384 & 2543, 4081, 2833\\
Square & With product & 16,384 & 16,384 & 16,384 & 1, 1, 1\\
Sine & Affine/ReLU & 8,191 & 9,272 & 5,662 & 4067, 3959, 3790\\
Sine & With product & 16,384 & 16,384 & 16,384 & 587, 526, 772\\
Positive self-power & Affine/ReLU & 0 & invalid & invalid & 4026, invalid, invalid\\
Positive self-power & With product & 16,384 & 16,384 & 16,380 & 2622, 3008, 2881\\
\bottomrule
\end{tabular}
\end{table}

On these same outer sets the three affine/rectifier sine circuits have raw RMSE $1.47\times10^{-16}$, $1.03\times10^{-16}$ and $1.07\times10^{-16}$. For affine/rectifier self-power, the sole valid circuit has RMSE 0.0166. With product gates, the third self-power circuit has RMSE $1.18\times10^{-13}$. All other evaluable rows and trials in Table~\ref{tab:circuit-synthesis} have RMSE zero. Each product-gate square reply uses a single binary64 multiplication, which computes the reference by definition.

Every affine/rectifier square proposal matched all 304 inputs specially chosen near rounding midpoints on the outer intervals and all 32,768 fresh ordinary outer inputs. Across 16,369 inputs sampled over a wider range of exponents, the three exact counts were 16,357, 16,294 and 16,364. Appendix~\ref{app:literal-square-certificate} proves the first circuit's result on both complete outer intervals. On 1,884 CORE-MATH sine inputs in those intervals the three affine/rectifier replies matched 1,026, 1,079 and 697, compared with 1,881, 1,884 and 1,883 with product gates. The product-gate self-power circuits missed one, one and two of the 896 near-midpoint probes below one defined in Appendix~\ref{app:smooth-circuits}. On the 8,192-input wide-exponent panel, the second and third product-gate self-power replies each produced nonfinite outputs, so neither has finite raw RMSE there. The electronic results retain each trial, split, exact count and RMSE status. In this experiment Astra chose each topology. The same-graph BP comparisons in Appendix~\ref{app:smooth-optimizer} concern the three separately engineered smooth circuits.
